\documentclass[conference]{IEEEtran}

% ---- Packages ----
\usepackage{cite}
\usepackage{amsmath,amssymb,amsfonts}
\usepackage{graphicx}
\usepackage{textcomp}
\usepackage{xcolor}
\usepackage{booktabs}
\usepackage{multirow}
\usepackage{url}
\usepackage{hyperref}
\usepackage{balance}
\usepackage{enumitem}
\usepackage{array}
\usepackage{tabularx}

\hypersetup{
  colorlinks=true,
  linkcolor=black,
  citecolor=black,
  urlcolor=blue,
  pdfauthor={Arun Narayanan},
  pdftitle={Empirical Evaluation of Hallucinations in Large Language Models}
}

% Column type for tables
\newcolumntype{C}[1]{>{\centering\arraybackslash}p{#1}}
\newcolumntype{R}[1]{>{\raggedleft\arraybackslash}p{#1}}

\def\BibTeX{{\rm B\kern-.05em{\sc i\kern-.025em b}\kern-.08em
    T\kern-.1667em\lower.7ex\hbox{E}\kern-.125emX}}

\begin{document}

\title{Empirical Evaluation of Hallucinations in Large Language Models: A Comparative Study of Code Generation and Technical Knowledge Tasks}

\author{
  \IEEEauthorblockN{Arun Narayanan}
  \IEEEauthorblockA{
    School of Computing\\
    Dublin City University\\
    Dublin, Ireland\\
    arun.narayanan2@mail.dcu.ie
  }
}

\maketitle

% ============================================================
% ABSTRACT
% ============================================================
\begin{abstract}
Large Language Models (LLMs) are increasingly deployed in software development and technical workflows, yet their tendency to generate confident but incorrect outputs---hallucinations---remains a critical concern. This study presents a comprehensive empirical evaluation of hallucination rates in two state-of-the-art LLMs: Google Gemini~2.5 Flash and Cohere Command~R7B. We employ a dual-mode testing framework combining objective code execution testing (50~questions) with knowledge-based questions (50~questions) across eight technical domains. Crucially, we distinguish between \emph{hallucinations} (factually incorrect answers) and \emph{format violations} (correct answers embedded in extra text)---a distinction often conflated in prior work. Gemini achieved 89\% overall accuracy while Cohere achieved 65\%, a 24~percentage-point gap. Code execution tasks proved substantially easier (90\% average) than specialized technical knowledge tasks (63\% average). Performance degraded markedly for specialized domains, particularly Snowflake (40\% average) and Spark (55\% average). These findings reveal critical gaps in specialized technical knowledge with significant implications for software development and data engineering workflows.
\end{abstract}

\begin{IEEEkeywords}
Large Language Models, Hallucinations, Code Generation, Technical Knowledge, Objective Evaluation, Format Violations
\end{IEEEkeywords}

% ============================================================
% 1. INTRODUCTION
% ============================================================
\section{Introduction}

Large Language Models such as GPT-4~\cite{achiam2023gpt4}, Gemini, and Cohere are now routinely employed for tasks ranging from general writing to specialized domains including software development, data engineering, and database management. A well-documented limitation is their propensity to generate hallucinations---confident, plausible-sounding outputs that are factually incorrect or logically flawed~\cite{ji2023survey}. This presents serious risks in technical domains where incorrect code or misinformation can lead to security vulnerabilities, system failures, and data loss.

Most existing evaluations of LLM performance rely on subjective human judgment or multiple-choice assessments, both of which suffer from limitations. In contrast, code generation and technical knowledge retrieval offer opportunities for objective, reproducible evaluation. When an LLM generates Python code, that code can be executed in a controlled environment and compared against expected outputs with binary precision.

This research addresses three specific questions:
\begin{enumerate}[leftmargin=*]
  \item How do state-of-the-art LLMs perform on objective coding tasks with clear ground truth?
  \item Do hallucination rates differ between code generation and factual knowledge retrieval tasks?
  \item Which technical domains are most prone to hallucinations, and what patterns emerge?
\end{enumerate}

We developed a comprehensive testing framework incorporating 100~questions spanning eight technical categories. The framework uses two evaluation modes: Python code execution via the Piston API sandbox for objective output verification, and direct string comparison for knowledge-based questions.

A critical contribution of this work is the explicit distinction between \textbf{hallucinations} (factually incorrect or logically flawed outputs) and \textbf{format violations} (correct answers embedded in unnecessary explanations or preambles). Prior work frequently conflates these categories, leading to overestimation of hallucination rates when models provide correct information with verbose formatting.

% ============================================================
% 2. METHODOLOGY
% ============================================================
\section{Methodology}

\subsection{Testing Framework Design}

\textbf{Mode~1: Python Code Execution (Objective Scoring).}
We prompt LLMs to generate complete Python code solving specific algorithmic and logical challenges. The generated code is transmitted to the Piston API (a cloud-based sandbox) where it is compiled and executed with a 5-second timeout. A response is marked correct only if: (a)~the code compiles without syntax errors, (b)~the code executes without runtime errors, and (c)~the output exactly matches the expected result character-for-character.

\textbf{Mode~2: Knowledge-Based Questions (String Matching with Format Detection).}
We prompt LLMs to answer technical questions with brief, precise answers. The answer is compared against ground truth using: (1)~exact matching, (2)~partial matching fallback, and (3)~format violation detection. We distinguish three response categories:
\begin{itemize}[leftmargin=*]
  \item \textbf{Correct}: Answer exactly matches expected output.
  \item \textbf{Format Violation}: Correct answer present but with extra text (e.g., ``The answer is DISTINCT'' instead of ``DISTINCT'').
  \item \textbf{Hallucination}: Factually incorrect or logically flawed answer.
\end{itemize}

System prompts explicitly instruct models to provide only answers without explanations, allowing us to quantify format compliance violations separately.

\subsection{Dataset Composition}

The dataset consists of \textbf{100~questions} divided equally between two modes.

\textit{Python Code Execution Questions (50~total):}
\textbf{Edge Cases}~(15): floating-point arithmetic, type comparison, identity vs.\ equality, truthiness evaluation.
\textbf{String Tricks}~(10): reversal, character counting, whitespace handling, palindrome detection.
\textbf{Algorithms}~(15): prime detection, GCD/LCM, factorial computation, perfect squares, divisor sums.
\textbf{Data Structures}~(10): list slicing, dictionary operations, set operations, matrix transposition.
Approximately 60\% are classified as \emph{Hard} difficulty, targeting edge cases to maximize discrimination between models.

\textit{Knowledge-Based Questions (50~total):}
\textbf{Polars}~(15): modern DataFrame library covering lazy evaluation (\texttt{scan\_csv}, \texttt{collect}), column operations, conditional operations, joins, and window functions.
\textbf{SQL}~(15): standard SQL including \texttt{SELECT}, \texttt{JOIN}, \texttt{GROUP BY}, aggregate functions, and set operations.
\textbf{Spark}~(10): distributed computing concepts including actions, transformations, persistence levels, and shuffle operations.
\textbf{Snowflake}~(10): cloud warehouse features including \texttt{QUALIFY}, \texttt{MAX\_BY}, \texttt{GROUP BY ALL}, \texttt{EXCLUDE}, Time Travel, and \texttt{FLATTEN}.

The dataset intentionally includes specialized, modern tools (Polars, Snowflake) where training data may be limited, enabling evaluation of model behavior at the frontier of their knowledge.

\subsection{AI Models Tested}

\textbf{Google Gemini~2.5 Flash}: Released July~2025, lightweight multimodal transformer with a 1-million-token context window, accessed via the Google Gemini API.

\textbf{Cohere Command~R7B}: Released December~2024, 7-billion-parameter decoder-only model designed for enterprise deployments, accessed via the Cohere~v2 Chat API (model~identifier: \texttt{command-r7b-12-2024}).

Both models were evaluated through official commercial APIs with identical prompts and evaluation criteria. No fine-tuning or model-specific optimizations were applied.

\subsection{Evaluation Criteria}

\textit{Python Code Responses:} Syntax correctness, runtime correctness, output correctness, and timeout handling (5-second limit). Correct if and only if all criteria pass.

\textit{Knowledge Question Responses:} Exact match $\rightarrow$ Correct~(1.0); partial match without explanation $\rightarrow$ Correct~(1.0); correct answer + explanation $\rightarrow$ Format Violation~(0.5); incorrect answer $\rightarrow$ Hallucination~(0.0). Format violations and hallucinations are tracked separately.

\subsection{Prompt Engineering}

Python mode: \emph{``Write complete Python code that solves this problem. Only provide the code, no explanations.''}
Knowledge mode: \emph{``Answer with ONLY the exact function name, keyword, or term. No explanations, no extra text.''}

Both prompts were deliberately straightforward, designed to test models on their default behavior rather than highly engineered responses.

\subsection{Statistical Methods}

All reported accuracy percentages include 95\% confidence intervals calculated using the exact binomial distribution. Category-level differences are tested using Fisher's exact test (for groups with $n<20$) with Bonferroni correction applied for multiple comparisons ($\alpha = 0.05/8 = 0.006$ per category). Effect sizes are reported using Cram\'{e}r's~$V$.

% ============================================================
% 3. RESULTS
% ============================================================
\section{Results}

\subsection{Overall Performance}

Gemini~2.5 Flash achieved \textbf{89\% overall accuracy} (89/100), while Cohere Command~R7B achieved \textbf{65\% overall accuracy} (65/100), a \textbf{24~percentage-point gap}. The 95\% confidence intervals do not overlap (Gemini:~[81--95\%]; Cohere:~[55--74\%]), confirming the gap is statistically significant.

\begin{table}[!t]
\caption{Overall Performance Summary}
\label{tab:overall}
\centering
\begin{tabular}{l r r r}
\toprule
\textbf{Metric} & \textbf{Gemini} & \textbf{Cohere} & \textbf{$\Delta$} \\
\midrule
Overall Accuracy       & 89\% & 65\% & +24\,pp \\
95\% CI                & [81--95\%] & [55--74\%] & --- \\
Correct Responses      & 89  & 65  & +24 \\
Format Violations      & 6   & 19  & $-$13 \\
Hallucinations         & 5   & 16  & $-$11 \\
\bottomrule
\end{tabular}
\vspace{2pt}
\footnotesize Format violations and hallucinations tracked separately. Cohere's higher format violation rate (19\% vs.\ 6\%) accounts for much of the gap.
\end{table}

\subsection{Performance by Mode}

Performance diverges sharply between code execution and knowledge tasks (Table~\ref{tab:mode}).

\begin{table}[!t]
\caption{Performance by Evaluation Mode}
\label{tab:mode}
\centering
\begin{tabular}{l r r r}
\toprule
\textbf{Mode} & \textbf{Gemini} & \textbf{Cohere} & \textbf{Avg} \\
\midrule
Python Code    & 96\% (48/50) & 84\% (42/50) & 90\% \\
Knowledge      & 82\% (41/50) & 46\% (23/50) & 64\% \\
\midrule
\textbf{Overall} & \textbf{89\%} & \textbf{65\%} & \textbf{77\%} \\
\bottomrule
\end{tabular}
\vspace{2pt}
\footnotesize The 26\,pp gap between Python (90\%) and Knowledge (64\%) average accuracy indicates that objective, syntax-constrained tasks are significantly easier for LLMs.
\end{table}

The 26~percentage-point difference between Python (90\% average) and Knowledge (64\% average) demonstrates that objective, rule-based tasks with clear syntax constraints are significantly easier for LLMs than open-ended knowledge retrieval. Format violations are more prevalent in knowledge tasks (Cohere: 28\% in knowledge vs.\ 10\% in code).

\subsection{Performance by Category}

Tables~\ref{tab:python} and~\ref{tab:knowledge} present category-level results with statistical significance annotations.

\begin{table}[!t]
\caption{Python Code Category Performance}
\label{tab:python}
\centering
\begin{tabular}{l r r r r}
\toprule
\textbf{Category} & \textbf{Gem.} & \textbf{Coh.} & \textbf{Avg} & \textbf{$\Delta$} \\
\midrule
Edge Case (15)      & 93\% & 80\% & 87\% & +13\,pp \\
String Trick (10)   & 100\% & 90\% & 95\% & +10\,pp \\
Algorithm (15)      & 100\% & 87\% & 93\% & +13\,pp \\
Data Structure (10) & 90\% & 78\% & 84\% & +12\,pp \\
\midrule
\textbf{Average}    & \textbf{96\%} & \textbf{84\%} & \textbf{90\%} & \textbf{+12\,pp} \\
\bottomrule
\end{tabular}
\vspace{2pt}
\footnotesize None of the Python category differences reach significance after Bonferroni correction ($\alpha=0.006$), reflecting relatively similar performance on code tasks.
\end{table}

\begin{table}[!t]
\caption{Knowledge Domain Category Performance}
\label{tab:knowledge}
\centering
\begin{tabular}{l r r r r l}
\toprule
\textbf{Category} & \textbf{Gem.} & \textbf{Coh.} & \textbf{Avg} & \textbf{$\Delta$} & \textbf{Sig.} \\
\midrule
Polars (15)     & 87\% & 53\% & 70\% & +34\,pp & * \\
SQL (15)        & 87\% & 80\% & 83\% & +7\,pp  & ns \\
Spark (10)      & 80\% & 30\% & 55\% & +50\,pp & ** \\
Snowflake (10)  & 70\% & 10\% & 40\% & +60\,pp & *** \\
\midrule
\textbf{Average} & \textbf{82\%} & \textbf{46\%} & \textbf{64\%} & \textbf{+36\,pp} & \\
\bottomrule
\end{tabular}
\vspace{2pt}
\footnotesize Significance (Fisher's exact, Bonferroni-corrected): ns = not significant; * $p<0.05$; ** $p<0.01$; *** $p<0.001$.
\end{table}

The large gaps in specialized domains (Snowflake: 60\,pp, Spark: 50\,pp) are statistically robust, while differences in fundamental domains (SQL: 7\,pp) are not significant.

\subsection{Format Violations vs.\ Hallucinations}

Approximately 8--15\% of responses classified as failures are \emph{format violations} (correct answer with extra text) rather than hallucinations. Three illustrative examples:

\textbf{Example~1} (Knowledge). Question: ``What SQL keyword retrieves unique values?'' Expected: \texttt{DISTINCT}. Cohere: ``SELECT DISTINCT to retrieve unique values from a column.'' Classification: Format Violation.

\textbf{Example~2} (Knowledge). Question: ``What Polars method scans a CSV file lazily?'' Expected: \texttt{scan\_csv}. Cohere: ``Polars offers the scan\_csv method for lazy loading of CSV files.'' Classification: Format Violation.

\textbf{Example~3} (Knowledge). Question: ``What SQL function returns the maximum value?'' Expected: \texttt{MAX}. Cohere: ``The MAX function returns the maximum value.'' Classification: Format Violation.

After reclassification: Gemini shows 6\% format violations and 5\% hallucinations; Cohere shows 19\% format violations (predominantly in knowledge tasks) and 16\% hallucinations.

\subsection{Failure Mode Analysis}

\textbf{Gemini's failures} (11~total: 6~format violations, 5~hallucinations): Hallucinations include incorrect identity comparisons, wrong modulo results, and confused function names (\texttt{filter} vs.\ \texttt{select} in Polars).

\textbf{Cohere's failures} (35~total: 19~format violations, 16~hallucinations): Hallucinations include plausible but incorrect code, confused function names in Spark/Snowflake, and incorrect terminology. Format violations are predominantly in knowledge questions where correct answers are wrapped in explanations.

% ============================================================
% 4. DISCUSSION
% ============================================================
\section{Discussion}

\subsection{Why Knowledge Tasks Are Harder}

Performance on knowledge tasks (64\% average) lags code execution (90\% average) by 26~percentage points. Four factors contribute:

\textbf{Domain Specialization \& Training Data.} Polars (released~$\sim$2020) and modern Snowflake features (2023--2024) post-date most LLM training cutoffs. Knowledge questions were intentionally selected from modern, specialized domains where models are expected to have gaps.

\textbf{Terminology Precision.} Knowledge answers require exact function or keyword names. LLMs frequently hallucinate plausible alternatives---e.g., \texttt{read\_csv} instead of \texttt{scan\_csv} (Polars), \texttt{WHERE} instead of \texttt{QUALIFY} (Snowflake). Each alternative is syntactically plausible but incorrect.

\textbf{Absence of Syntax Validation.} Code execution has a built-in validation mechanism: syntax and runtime errors are caught immediately. Knowledge questions have no such compilation phase---incorrect terminology passes silently.

\textbf{Disambiguation Requirements.} Knowledge questions reference functions within large multi-function libraries (Polars has 300+ methods). The LLM must disambiguate among semantically similar options, whereas code questions typically have a unique correct solution.

\subsection{Specialized Domains: Snowflake at 40\%}

Snowflake achieved the worst category performance (Gemini~70\%, Cohere~10\%, average~40\%). Root causes include: (1)~enterprise tool specificity---Snowflake is closed-source with documentation underrepresented in public training data; (2)~recent feature additions (\texttt{QUALIFY}, \texttt{MAX\_BY}) introduced in 2023--2024, potentially after training cutoffs; (3)~unique syntax patterns diverging from standard SQL; and (4)~Cohere's smaller model capacity (7B~parameters) limiting knowledge-intensive recall.

\subsection{Cohere's Format Violation Tendency}

Cohere demonstrated a 19\% format violation rate compared to Gemini's 6\%. Despite explicit instructions to provide only answers, Cohere frequently appended explanations. This suggests Cohere's training emphasizes helpful, explanatory responses, making it less effective at suppressing verbosity. Importantly, format violations indicate that \emph{knowledge is present} but instruction compliance is lacking---a fundamentally different failure mode from hallucinations, and one potentially recoverable through output post-processing.

% ============================================================
% 5. THREATS TO VALIDITY
% ============================================================
\section{Threats to Validity}

\subsection{Construct Validity}

Our 100~test questions focus on structured, objectively-scored tasks. Open-ended generation is not evaluated. Python questions emphasize syntax and logic, not algorithmic efficiency. Knowledge questions test terminology recall, not conceptual understanding. Results represent hallucination rates for \emph{this specific evaluation domain}, not all possible LLM applications.

\subsection{Internal Validity}

We use a single prompt formulation per category; response patterns might change with different prompts. All generations used default temperature settings. We evaluate zero-shot performance only---few-shot examples might improve results. The Piston API sandbox's 5-second timeout may not reflect all production scenarios.

\subsection{External Validity}

Only two models were evaluated. Results do not generalize to GPT-4, Claude~3, Llama~3, or other LLMs. Model size disparity (Cohere~7B vs.\ Gemini's unknown but likely larger architecture) affects comparisons. Our knowledge questions focus on modern data tools; generalization to medical, legal, or scientific domains is unclear. All questions are English-only.

\subsection{Conclusion Validity}

Binary scoring misses partial correctness. Category sample sizes are small ($n=10$--$15$), yielding wide confidence intervals. We applied Bonferroni correction to control false positives across 8~category comparisons. No inter-rater reliability study was conducted---our exact-match criterion is stricter than typical human judgment.

% ============================================================
% 6. RELATED WORK
% ============================================================
\section{Related Work}

\subsection{Hallucination Definitions}

Ji et~al.~\cite{ji2023survey} provide a comprehensive survey distinguishing extrinsic hallucinations (contradicting external knowledge) from intrinsic hallucinations (self-contradictory). Our work focuses on extrinsic hallucinations. Maynez et~al.~\cite{maynez2020faithfulness} note that ``correct facts delivered in verbose formatting'' are fundamentally different from factual errors---a distinction we operationalize as format violations vs.\ hallucinations.

\subsection{Code Evaluation Benchmarks}

Chen et~al.~\cite{chen2021evaluating} present HumanEval, establishing execution-based code evaluation. Liu et~al.~\cite{liu2023humaneval} extend this with HumanEval+, adding robustness through diverse test cases. Our Python evaluation follows the HumanEval paradigm but focuses on edge cases, evaluates commercial LLMs, and distinguishes format violations from correctness failures.

\subsection{Knowledge \& Truthfulness Benchmarks}

Hendrycks et~al.~\cite{hendrycks2020measuring} introduce MMLU, testing knowledge across 57~domains via multiple-choice questions. Lin et~al.~\cite{lin2022truthfulqa} present TruthfulQA, evaluating factual alignment with open-ended answers. Min et~al.~\cite{min2023factscore} introduce FactScore for fine-grained evaluation of factual precision. Our distinction between format violations and hallucinations echoes FactScore's attention to granular evaluation.

\subsection{Domain-Specific LLM Evaluation}

Rozi\`{e}re et~al.~\cite{roziere2024code} present Code~Llama, evaluated on programming benchmarks. Dashwood~\cite{dashwood2023cohere} compares Cohere models across diverse tasks. Our work uniquely evaluates \emph{modern data engineering domains} (Snowflake, Polars, Spark) underrepresented in prior benchmarks.

\subsection{Positioning of This Work}

This research contributes: (1)~explicit format violation detection, separating verbose correctness from hallucinations; (2)~specialized technical domain evaluation extending beyond foundational domains; (3)~comparative binary evaluation demonstrating that code tasks are substantially easier than knowledge retrieval; and (4)~clear, reproducible methodology compatible with HumanEval~\cite{chen2021evaluating} and TruthfulQA~\cite{lin2022truthfulqa}.

% ============================================================
% 7. LIMITATIONS & FUTURE WORK
% ============================================================
\section{Limitations \& Future Work}

\subsection{Limitations}

\textbf{Prompt Sensitivity.} We use single, straightforward prompts. More aggressive prompting or few-shot examples might improve results; intentionally poor prompts could degrade them.

\textbf{Model Coverage.} Only two models were evaluated. Extending to GPT-4, Claude~3, Llama~3 would strengthen conclusions.

\textbf{Piston API Deprecation.} The Piston public API transitioned to whitelist-only access in February~2026. This evaluation was completed before the transition. Reproducibility requires self-hosted Piston instances or alternative sandboxes. We provide local execution scripts in our GitHub repository~\cite{repo}.

\textbf{Dataset Authorship.} Test questions were created by the author without independent expert verification. Domain expert review would strengthen validity.

\textbf{Training Data Cutoffs.} We lack access to exact training data cutoffs, limiting our ability to distinguish training data gaps from architectural limitations.

\subsection{Future Work}

Key directions include: (1)~expanded model evaluation across GPT-4, Claude~3, Llama~3, and open-source models; (2)~few-shot learning with 1, 3, and 5 in-context examples; (3)~systematic prompt variation studies; (4)~temperature and sampling parameter sensitivity analysis; (5)~human evaluation with inter-rater reliability for knowledge questions; and (6)~larger context evaluation to assess context-dependence.

% ============================================================
% 8. CONCLUSION
% ============================================================
\section{Conclusion}

This empirical evaluation reveals clear hierarchies in LLM performance across technical domains.

\textbf{Key Findings:}
(1)~Gemini~2.5 Flash substantially outperforms Cohere Command~R7B (89\% vs.\ 65\%), with the largest gaps in specialized domains (Snowflake: 60\,pp gap).
(2)~Code execution tasks are substantially easier than knowledge retrieval (90\% vs.\ 64\% average), suggesting that objective constraints and syntax validation strongly improve LLM performance.
(3)~Format violations account for 8--15\% of failures---these are correct answers with verbose formatting, distinct from true hallucinations.
(4)~Specialized tools (Snowflake, Polars) expose significant gaps, with performance degrading to 40\% in the worst case.

\textbf{Practical Implications:}
For critical applications (production code, database queries), human review remains essential---neither model achieves near-perfect accuracy across all domains. Format violations suggest that output post-processing might recover some responses; true hallucinations are not recoverable. Practitioners should not assume LLMs have reliable knowledge of modern specialized tools without domain-specific validation.

\textbf{Methodological Contributions:}
This work demonstrates the value of objective, executable evaluation methodologies with binary, reproducible scoring. The explicit distinction between format violations and hallucinations provides a more nuanced understanding of LLM failure modes. All code and data are publicly available.\footnote{\url{https://github.com/ArunGMR0411/ai-hallucination-eval}}

\balance
\bibliographystyle{IEEEtran}
\bibliography{references}

\end{document}
